\section[Removal of restriction on $\arg z$]{Removal of restriction on $\boldsymbol{\arg z}$}\label{removal}

Using $\beta(u)$ from Lemma \ref{2:l2} we def\/ine
\begin{gather*} W_4(u,z)=\frac{\pi i}{\beta(u)} W_2(u,z) .\end{gather*}
Then we have
\begin{gather}\label{3:connect1}
W_4\big(u,ze^{\pi i}\big)=e^{\pi i(1-\mu)} W_4(u,z)+\pi i W_3(u,z) .
\end{gather}
Moreover, \eqref{2:beta} shows that $W_4$ shares the asymptotic expansion \eqref{1:W2}, \eqref{1:W2a} with~$W_2$.
From~\eqref{3:connect1} we obtain
\begin{gather}\label{3:connect}
W_4\big(u,ze^{\pi im}\big)=e^{\pi i (1-\mu)m} W_4(u,z)+\pi i\frac{\sin(\pi(\mu+1)m)}{\sin(\pi(\mu+1))} W_3(u,z)
\end{gather}
for every integer $m$.
We will use~\eqref{3:connect} and the asymptotic expansions \eqref{1:W1}, \eqref{1:W2} for $|\arg z|\le \frac12\pi$
to prove that in~\eqref{1:W2a} we can remove the restriction on~$\arg z$ completely.

\begin{Theorem}\label{3:t1}
Suppose that $\Re\mu\ge0$. For every $N=1,2,3,\dots$, $W_2(u,z)$ can be written as the right-hand side of~\eqref{1:W2},
and \eqref{1:W2a} holds without a restriction on~$\arg z$:
\begin{gather*}%\label{3:W2b}
|g_2(u,z)|+|h_2(u,z)|\le \frac{K_2}{u^{2N}}\qquad\text{for} \quad 0<|z|\le R, \quad u\ge u_2.
\end{gather*}
\end{Theorem}

\begin{proof}
Without loss of generality we replace $W_2$ by $W_4$. We assume that $|\arg z|\le \frac12\pi$, $0<|z|\le R$, $u>0$, $m$ is an integer and $z_1:=ze^{\pi im}$.
We insert~\eqref{1:W1}, \eqref{1:W2} on the right-hand side of~\eqref{3:connect}.
Using~\eqref{3:ancontK} we obtain
\begin{gather}\label{3:eq4}
W_4(u,z_1)=z_1K_\mu(uz_1)\sum_{s=0}^{N-1} \frac{A_s(z_1)}{u^{2s}}
-\frac{z_1}{u}K_{\mu+1}(uz_1)\sum_{s=0}^{N-1} \frac{B_s(z_1)}{u^{2s}}+f(u,z),
\end{gather}
where
\begin{gather*}%\label{3:f}
 f=E_1g_2+E_2h_2+E_3g_1+E_4h_1,
\end{gather*}
with
\begin{alignat*}{3}
& E_1(u,z) = e^{-\pi i(\mu+1)m} zK_\mu(uz),\qquad &&
 E_2(u,z) = -e^{-\pi i (\mu+1)m}\frac{z^2}{u} K_{\mu+1}(uz),&\\
& E_3(u,z) = \pi i\frac{\sin(\pi(\mu+1)m)}{\sin(\pi(\mu+1))}z I_\mu(uz),\qquad&&
 E_4(u,z) = \pi i \frac{\sin(\pi (\mu+1)m)}{\sin(\pi (\mu+1))}\frac{z^2}{u}I_{\mu+1}(uz) .&
\end{alignat*}
We will construct functions $G_j(u,z)$ and $H_j(u,z)$ such that
\begin{gather*} E_j(u,z)=z_1 K_\mu(uz_1) G_j(u,z)-\frac{z_1^2}{u} K_{\mu+1}(uz_1) H_j(u,z)
\end{gather*}
for $j=1,2,3,4$.
Then \eqref{3:eq4} becomes
\begin{gather}
W_4(u,z_1)=z_1K_\mu(uz_1)\left(\sum_{s=0}^{N-1} \frac{A_s(z_1)}{u^{2s}}+g_3(u,z)\right)\nonumber\\
\hphantom{W_4(u,z_1)=}{}
-\frac{z_1}{u}K_{\mu+1}(uz_1)\left(\sum_{s=0}^{N-1} \frac{B_s(z_1)}{u^{2s}}+z_1h_3(u,z)\right),\label{4:eq5}
\end{gather}
where
\begin{gather*}
g_3 = G_1 g_2+G_2h_2+G_3 g_1+G_4 h_1,\qquad
h_3 = H_1 g_2+H_2h_2+H_3 g_1+H_4 h_1.
\end{gather*}
We now use \cite[(10.28.2)]{NIST}
\begin{gather}\label{3:relation}
 K_\mu(x)I_{\mu+1}(x)+K_{\mu+1}(x)I_\mu(x)=\frac1x .
\end{gather}
From \eqref{3:relation} and the relation
\begin{gather}\label{3:ancontI}
 I_\mu\big(ze^{\pi im}\big)=e^{\pi i\mu m}I_\mu(z)
\end{gather}
we obtain
\begin{gather*} u z_1K_\mu(uz_1) e^{\pi i (\mu+1)m} I_{\mu+1}(uz)+uz_1K_{\mu+1}(uz_1) e^{\pi i \mu m} I_\mu(uz)=1 .\end{gather*}
Therefore, we can choose
\begin{gather*}
G_1(u,z) = uz K_\mu(uz)I_{\mu+1}(uz),\qquad
H_1(u,z) = -u^2K_\mu(uz)I_\mu(uz) .
\end{gather*}
We set
\begin{gather*}
l_0(x)=\ln\frac{1+2|x|}{|x|},\qquad l_\mu(x)=1\qquad\text{if} \quad \mu\ne 0,
\end{gather*}
and note the estimates \cite[(9.12)]{O}
\begin{alignat}{3}
&|I_\mu(x)K_\mu(x)|\le \frac{C l_\mu(x)}{1+|x|},\qquad&& |I_{\mu+1}(x)K_\mu(x)|\le \frac{C |x|l_\mu(x)}{1+|x|^2},& \label{3:olver1}\\
& |I_{\mu+1}(x)K_{\mu+1}(x)|\le \frac{C}{1+|x|},\qquad&& |I_\mu(x)K_{\mu+1}(x)|\le \frac{C}{|x|}& \label{3:olver2}
\end{alignat}
valid when $|\arg x|\le \frac12\pi$ with $C$ independent of $x$.
At this point we assume that $\mu\ne0$ (the case $\mu=0$ is mentioned at the end of the proof). The estimates~\eqref{3:olver1} give
\begin{gather}\label{3:est1}
|G_1(u,z)|\le C,\qquad |H_1(u,z)|\le C u^2.
\end{gather}
Similarly, we choose
\begin{gather*}
G_2(u,z) = -z^2K_{\mu+1}(uz)I_{\mu+1}(uz),\qquad
H_2(u,z) = uzK_{\mu+1}(uz)I_\mu(uz).
\end{gather*}
The estimates \eqref{3:olver2} give
\begin{gather}\label{3:est2}
|G_2(u,z)|\le C|z|^2,\qquad|H_2(u,z)|\le C.
\end{gather}
It follows from \eqref{3:ancontK} that
\begin{gather*}
E_3(u,z) = -E_1(u,z)+ z_1K_\mu(uz_1),\qquad
E_4(u,z) = -E_2(u,z)-\frac{z_1^2}{u}K_{\mu+1}(uz_1) .
\end{gather*}
Therefore, we can choose
\begin{alignat*}{3}
& G_3(u,z) = 1-G_1(u,z) ,\qquad &&
H_3(u,z) = -H_1(u,z),&\\
& G_4(u,z) = -G_2(u,z),\qquad&&
H_4(u,z) = 1-H_2(u,z).&
\end{alignat*}
From \eqref{3:est1}, \eqref{3:est2}, we get
\begin{alignat}{3}
& |G_3(u,z)|\le C+1,\qquad && |H_3(u,z)|\le C u^2,& \label{3:est3}\\
& |G_4(u,z)|\le C|z|^2,\qquad && |H_4(u,z)|\le C+1 .&\label{3:est4}
\end{alignat}
The estimates \eqref{3:est1}, \eqref{3:est2}, \eqref{3:est3}, \eqref{3:est4}
give
\begin{gather*}
|g_3(u,z)| \le C |g_2(u,z)| + C|z|^2|h_2(u,z)|+(C+1)|g_1(u,z)|+C|z|^2 |h_1(u,z)|,\\
|h_3(u,z)| \le C u^2|g_2(u,z)| + C|h_2(u,z)|+Cu^2|g_1(u,z)|+(C+1) |h_1(u,z)|.
\end{gather*}
Since we assumed that
\begin{gather*}
|g_1(u,z)|+|h_1(u,z)|+|g_2(u,z)|+|h_2(u,z)|\le \frac{K}{u^{2N}}
\end{gather*}
for $| \arg z|\le \frac12\pi$, $0<|z|\le R$, $u\ge u_0$,
the expansion \eqref{4:eq5} has the desired form with $N$ replaced by $N-1$.

Suppose $\mu=0$.
We use \cite[(10.31.2)]{NIST}
\begin{gather}\label{3:K0}
K_0(x)=-\left(\ln\left(\frac12x\right)+\gamma\right)I_0(x)+\frac{\frac14x^2}{(1!)^2}+\left(1+\frac12\right)
\frac{\left(\frac14x^2\right)^2}{\left(2!\right)^2}+\cdots.
\end{gather}
It follows from \eqref{3:K0} that there exist positive constants $r>0$, $D>0$ such that
\begin{gather*}%\label{3:est5}
\frac{|K_0(x)|}{|K_0(xe^{\pi i m})|}\le D \qquad\text{for} \quad 0<|x|\le r, \quad |\arg x|\le\frac12\pi, \quad m\in\mathbb{Z}.
\end{gather*}
Then we set
\begin{gather*}
G_1(u,z)=\frac{K_0(uz)}{K_0(uz_1)},\qquad H_1(u,z)=0\qquad\text{if} \quad 0<|uz|\le r
\end{gather*}
with $G_1$ and $H_1$ the same as before when $|uz|>r$. The estimates~\eqref{3:est1} are valid with a suitable constant~$C$.
The rest of the proof is unchanged.
This completes the proof of the theorem.
\end{proof}
